% TOP_PROBLEM: {"id": 30002442, "title": "Homogeneity of Fano Contact Manifolds", "category_id": 6, "category": {"id": 6, "name": "geometry", "display_name": "Geometry", "description": "Euclidean and non-Euclidean geometry, geometric structures.", "slug": "geometry", "order_index": 6, "created_at": "2026-07-31T15:26:25.671Z"}, "status": "open"}
% FIRST_POSTED: 2026-10-06
% ATTEMPT_STATUS: solved
% OPENAI_CLAIM: {"schema_version":1,"source_repo":"https://github.com/openai/math","source_commit":"adc7f1241b42e322a6451854ab7e4b4c146bf78a","release_date":"2026-10-06","match":"full","resolution":"proved","independently_reviewed":false,"second_pass":"agreed","families":[{"family":"062","title":"Projective contact classification and the LeBrun–Salamon conjecture","subject":"Algebraic and complex geometry","manuscripts":[{"title":"Contact Fano manifolds and the LeBrun–Salamon conjecture","pdf_path":"preprints/Contact-Fano-manifolds-and-the-LeBrun-Salamon-conjecture-September-23-2026/paper.pdf","readme_path":"preprints/Contact-Fano-manifolds-and-the-LeBrun-Salamon-conjecture-September-23-2026/README.md","date":"2026-09-23","theorem_ref":"thm:main (sections/introduction.tex:30); thm:quaternionic (sections/introduction.tex:105)","statement_sources":[{"path":"preprints/Contact-Fano-manifolds-and-the-LeBrun-Salamon-conjecture-September-23-2026/build/sections/introduction.tex","line":30,"label":"thm:main"},{"path":"preprints/Contact-Fano-manifolds-and-the-LeBrun-Salamon-conjecture-September-23-2026/build/sections/introduction.tex","line":105,"label":"thm:quaternionic"}]}],"lean":null,"reasoning_trace":null}]}
\documentclass[11pt]{article}
\usepackage[margin=25mm]{geometry}
\usepackage{fontspec}
\setmainfont{Latin Modern Roman}
\usepackage{amsmath,amssymb,mathtools}
\usepackage{unicode-math}
\setmathfont{Latin Modern Math}
\usepackage{enumitem}
\usepackage{xurl,hyperref,bookmark}
\hypersetup{colorlinks=true,urlcolor=blue}
\setcounter{secnumdepth}{0}
\setlength{\parindent}{0pt}
\setlength{\parskip}{5pt}
\setlength{\emergencystretch}{3em}
\begin{document}
\section*{Homogeneity of Fano Contact Manifolds}

\subsection{OpenAI's claim}
Every complex contact Fano manifold is claimed to be an adjoint variety with its given contact distribution; the quaternionic-Kähler consequence identifies positive manifolds of real dimension $4m$, $m\ge2$, with symmetric Wolf models up to scale. The cited manuscript argues for that result; its relation to this problem is specified below.

\subsection{Mathematical statement of the claimed result}
\paragraph{Mathematical claim 1.} For every smooth connected projective complex contact Fano manifold $(X,F)$ of dimension $2n+1$, $n\ge1$, with contact line $L=T_X/F$, OpenAI claims a simple complex Lie algebra $\mathfrak g$ and an algebraic contact isomorphism
\[
X\simeq\mathbb P(\mathcal O_{\min}(\mathfrak g)),\qquad
L\simeq\mathcal O_{\mathbb P(\mathfrak g)}(1)|_X.
\]
The target is the projectivized least-dimensional nonzero nilpotent adjoint orbit, with its canonical contact distribution. The Riemannian corollary says that every closed connected $4m$-manifold, $m\ge2$, with holonomy contained in $\operatorname{Sp}(m)\operatorname{Sp}(1)$ and positive scalar curvature has $\nabla\operatorname{Rm}=0$ and agrees up to constant metric scaling with a Wolf space, which is compact, symmetric and positive quaternionic-Kähler.

\subsection{Outline of the argument}
Manuscript 1 (family 062, 2026-09-23): Two contact lines with nonzero Levi pairing are smoothed into very free conics while tracking both endpoint velocities. This produces a symmetric line-valued section on the square with the prescribed diagonal jet, forcing global contact fields to span; homogeneity follows, and the twistor correspondence supplies the Riemannian conclusion.

\subsection{Scope relative to this problem}
The claim identifies every smooth projective contact Fano manifold of dimension at least three, with its contact distribution, with a homogeneous adjoint variety. This implies the recorded homogeneity assertion. In complex dimension one the Fano curve is $\mathbb P^1$, which is homogeneous as well. The additional contact-line identification is compatible with, and stronger than, homogeneity alone.

\subsection{Formal verification}
Family 062 has no Lean scope page in this snapshot; no formalization of these manuscript results is listed.

This site has not executed the Lean files or checked their dependencies. These descriptions report the scope of the linked declarations, rather than a new verification of the proofs.

\subsection{Status caveat}
OpenAI attributes these results to an internal model. Its release notice cautions that results lacking full formalization can contain mistakes. This site has not independently assessed the proofs; a solved label records the scope of the reported claim.

\subsection{Sources}
\begin{itemize}
\item OpenAI, \emph{Contact Fano manifolds and the LeBrun–Salamon conjecture}, 2026-09-23; citation key \texttt{OAI:Contact-Fano-manifolds-and-the-LeBrun-Salamon-conjecture-September-23-2026}. \href{https://github.com/openai/math/blob/adc7f1241b42e322a6451854ab7e4b4c146bf78a/preprints/Contact-Fano-manifolds-and-the-LeBrun-Salamon-conjecture-September-23-2026/paper.pdf}{PDF (GitHub)}; \href{https://github.com/openai/math/raw/adc7f1241b42e322a6451854ab7e4b4c146bf78a/preprints/Contact-Fano-manifolds-and-the-LeBrun-Salamon-conjecture-September-23-2026/paper.pdf}{PDF (direct)}; \href{https://github.com/openai/math/blob/adc7f1241b42e322a6451854ab7e4b4c146bf78a/preprints/Contact-Fano-manifolds-and-the-LeBrun-Salamon-conjecture-September-23-2026/README.md}{Manuscript page and citation}; \href{https://github.com/openai/math/blob/main/preprints/Contact-Fano-manifolds-and-the-LeBrun-Salamon-conjecture-September-23-2026/paper.pdf}{Latest version}; Statement source: \href{https://github.com/openai/math/blob/adc7f1241b42e322a6451854ab7e4b4c146bf78a/preprints/Contact-Fano-manifolds-and-the-LeBrun-Salamon-conjecture-September-23-2026/build/sections/introduction.tex\#L30}{thm:main}; Statement source: \href{https://github.com/openai/math/blob/adc7f1241b42e322a6451854ab7e4b4c146bf78a/preprints/Contact-Fano-manifolds-and-the-LeBrun-Salamon-conjecture-September-23-2026/build/sections/introduction.tex\#L105}{thm:quaternionic}.
\item \href{https://github.com/openai/math/blob/adc7f1241b42e322a6451854ab7e4b4c146bf78a/CONTENTS.md}{OpenAI catalogue entry}.
\end{itemize}
\end{document}
